% !TEX root = Project3/Rapport/Project3.tex
\subsection{Decision trees with gradient boosting}
\label{sec:theory-xgboost}
Decision trees partition the feature space into boxes and
fit a constant in each, building the model by 
a sequence of simple yes-or-no questions. A single tree 
has low bias and very high variance; the remedy is to build
many trees and combine them, with various \textit{ensemble methods}
- baggin, random forests, AdaBoost - and the one used in this project: 
\textit{gradient boosting}, particularly the \textit{XGBoost} library. 
On tabular data, these ensemble methods are known to be the most accuracte, 
and may even outperform neural networks.
\cite[p.~199]{LectureNotes}. 
\subsubsection{Basics of a tree}
Before explaining XGBoost it is neccesary to explain the basics
of a decision tree, and its hyperparameters.
The \textit{root node} contains all the observations.
Each \textit{internal node} tests one feature against a threshold
and sends the observation left or right - \textit{branch}.
Each terminal node carries a prediction or classification - \textit{leaf}.
A classification or prediction means walking from the root to a leaf.
% Geometrically, each test .. feature space.. hyperplane .. axis j.

\subsubsection{Regression trees}
A single decision tree can handle both regression
and classification problems, and they only differ in what a leaf
predicts and how splits are scored. For a regression the the prediction
is the mean response in the leaf; for a classification tree it is
the most occuring class among the training observations in that leaf.
Although this project will solve
a classification problem, the decision trees temselves are regression
trees. % This is so for gradient boosting.
Regression trees will be explained briefly, before XGBoost is explained,
in which the tree building is optimized.  
% I dont understand
A regression tree is built in two steps. First the predictor space
is split into J distinct and non-overlapping regions. Second, 
for every observation falling into region j we make the same prediction,
the mean of the response values of the training observations in region j.
\subsubsection{XGBoost}
Well well well

\subsubsection{Implementation: Method}
The gradient boosting method applied in this project is from the
library \textit{XGBoost}:
